% Generated draft: requires mathematical and source review.
\documentclass[11pt]{article}
\usepackage[margin=25mm]{geometry}
\usepackage{fontspec}
\setmainfont{Latin Modern Roman}
\usepackage{amsmath,amssymb,mathtools,unicode-math}
\setmathfont{Latin Modern Math}
\usepackage{enumitem,needspace,xurl,hyperref}
\setlength{\parindent}{0pt}
\setlength{\parskip}{5pt}
\setlength{\emergencystretch}{3em}
\newcommand{\N}{\mathbb{N}}
\newcommand{\Z}{\mathbb{Z}}
\newcommand{\Q}{\mathbb{Q}}
\newcommand{\R}{\mathbb{R}}
\newcommand{\C}{\mathbb{C}}
\newcommand{\F}{\mathbb{F}}
\newcommand{\A}{\mathbb{A}}
\newcommand{\PP}{\mathbb{P}}
\newcommand{\E}{\mathbb{E}}
\newcommand{\eps}{\varepsilon}
\newcommand{\dd}{\,\mathrm{d}}
\begin{document}
\section*{UnsolvedMath ID: 20000972}
\textbf{Source problem code:} catalogue entry
\subsection{Definitions and mathematical statement}
For a finite loopless graph \(G\), an acyclic orientation has no directed cycle. The equivalence generated by turning a source into a sink corresponds to a special chip-firing class of divisors. The graphic arrangement has hyperplanes \(x_u=x_v\) for edges \(uv\), and chambers indexed by acyclic orientations; passing to \(\R^{V(G)}/\Z^{V(G)}\) gives the toric graphic arrangement. Matt Macauley's source question asks for a comparable geometric interpretation of chip-firing equivalence for larger divisor classes. It does not specify which larger classes or what geometric equivalence is required.

The recorded comments propose multi-toppling/hereditary chip-firing (Andrea Sportiello), ask whether multi-toppling is abelian (James Propp), assert that the abelian property holds precisely in the hereditary case of firing clusters and subclusters (Sportiello), and refer to Spencer Backman [6] (Lionel Levine). The PDF introduction explicitly says these are recorder summaries, not necessarily verbatim quotations; no universal abelian iff theorem is inferred without its firing conventions.

The supplied aid's partial construction fixes a connected finite loopless multigraph with a sink \(s\). Let \(V^\circ=V\setminus\{s\}\), let \(L\) be its reduced Laplacian, and let \(\mathcal H\) be a hereditary cover of \(V^\circ\): its members are nonempty subsets, it is closed under taking nonempty subsets, and their union is \(V^\circ\). For \(U\in\mathcal H\) define \(c_U(v)\), for \(v\in U\), as the number of edges from \(v\) to \(V\setminus U\), including sink edges and multiplicity. A nonnegative integer configuration \(a\) may fire \(U\) when
\[
 a_v\ge c_U(v)\quad(v\in U),
 \qquad a\longmapsto a-L\mathbf1_U.
\]
It is stable when no such cluster is legal. Define the monomial ideal, over any coefficient field,
\[
 J_{\mathcal H}=\left(\prod_{v\in U}x_v^{c_U(v)}:U\in\mathcal H\right).
\]
The aid reports that the standard monomials \(x^a\notin J_{\mathcal H}\) are exactly stable configurations. They form a bounded staircase, since all singletons are allowed and impose coordinate bounds. It also reports
\[
 \operatorname{span}_{\Z}\{L\mathbf1_U:U\in\mathcal H\}
 =L\Z^{V^\circ}.
\]
Thus the finite algebraic quotient \(\Z^{V^\circ}/L\Z^{V^\circ}\) is independent of the hereditary rule, while the stable staircase and the recurrent representative section attributed to Backman retain that rule. Stable configurations are not thereby asserted to be unique in each class. Equality of firing lattices concerns algebraic equivalence and is not, by itself, a statement of legal reachability between arbitrary nonnegative configurations; the recurrent dynamics and hypotheses must be specified.

For the ordinary cycle \(C_m\), \(m\ge3\), the aid gives a degree-zero completion. Write
\[
 \operatorname{Pic}^0(C_m)=\operatorname{Div}^0(C_m)/\operatorname{im}\Delta,
 \qquad D_O(v)=\operatorname{indeg}_O(v)-1.
\]
Since the cycle has equally many edges and vertices, \(D_O\) has degree zero. The reported map from source-to-sink classes of acyclic orientations \(O\) to \([D_O]\) is a bijection onto all nonzero classes of \(\operatorname{Pic}^0(C_m)\simeq\Z/m\Z\). Adjoining the one missing identity class completes this cycle case. This does not supply a toric-arrangement analogue for every larger divisor class on every graph, so it remains a partial answer to the original geometric question.
\subsection{Short English statement}
The source asks to extend the toric-arrangement picture of source-to-sink equivalence beyond orientation divisors. The supplied aid separates a fixed Laplacian quotient from rule-dependent stable configurations and reports a one-class completion for cycles, while retaining the unresolved geometric and firing-rule conventions.
\subsection{Sources}
\begin{itemize}
\item Local UnsolvedMath registry, permanent ID \(20000972\), source code \texttt{catalogue entry}; \texttt{aim-combinatorics-notes.json}, \texttt{notes[96]}, tag \texttt{problem}.
\item AIM, \emph{Generalizations of chip-firing and the critical group}, Problem 7, following the workshop of 8--12 July 2013, source reference [1]: \url{https://aimath.org/pastworkshops/chipfiringproblems.pdf}. Statement status \texttt{exact}. The supplied verification preserves the introduction's warning that comments are recorder summaries, not necessarily verbatim quotations. No PDF was accessed for this draft.
\item Matt Macauley, main question; Andrea Sportiello, multi-toppling and hereditary comments; James Propp, abelian-property question; Lionel Levine, reference to Spencer Backman [6]. The full bibliographic entry for [6] is not supplied.
\item Research attempt 1: \texttt{partial\_result}, status \texttt{partially\_solved}, classification \texttt{PARTIAL-PROGRESS}, rigor label \texttt{proved}, novelty confidence low. Created and updated 14 August 2026. This machine-generated aid is not peer reviewed; independently verify before citation.
\item Ulam AI / UnsolvedMath Contributors; CC BY 4.0.
\end{itemize}
\end{document}
